

\begin{defn}[Implications/Conditional Statements]
    A \emph{conditional statement}, or {\em implication}, is a compound statement of the form {\em ``if \(P\), then \(Q\).''}  The first statement, \(P\), is called the {\em hypothesis} or {\em premise} and the second is the {\em conclusion}.  Conditionals can also be read {\em ``\(P\) implies \(Q\),''} {\em ``\(P\) only if \(Q\),''}, {\em ``\(Q\), if \(P\).''}  Notationally, we write ``if \(P\), then \(Q\)'' as \(P\rightarrow Q\).
\end{defn}

\setcounter{n}{1}
\begin{practice}[label=GRimages]{When are Implications False}
\setlength{\extrarowheight}{10pt}
    Looking at the following images, 
    \begin{center}
        \begin{tikzpicture}
            \foreach \p/\x in {GR/0,G/4,R/8,None/12}{
                \pic[scale=0.5,transform shape] at (\x,0) {\p};
                \node at (\x-1,2.5) {(\then)};
                \stepcounter{n};
            }
        \end{tikzpicture}
    \end{center}
    for which images is the statement 
        \begin{center}\emph{``If it is raining, then the grass is wet'' (\(R\rightarrow G\))}\end{center}
    false?\vspace{\baselineskip}\newline\hspace*{0pt}\hrulefill\vspace{\baselineskip}\newline
    Why is it false in those cases?\hrulefill\vspace{0.75\baselineskip}\newline  
    \hspace*{0pt}\hrulefill\vspace{\baselineskip}\newline
    In what situations must it then be true?\hrulefill\vspace{\baselineskip}\newline
    \hspace*{0pt}\hrulefill\vspace{\baselineskip}\newline
    Use the same sort of reasoning to try and fill in the following \emph{truth table}\index{truth table}.
    \begin{center}
    \begin{tabular}{*{4}{|c}*{5}{|>{\centering}m{0.125\textwidth}}|}\hline
        \(R\) & \(G\) & \(\sim R\) & \(\sim G\) & \(R\rightarrow G\) & \(G\rightarrow R\) & \(\sim R\rightarrow \sim G\) & \(\sim G\rightarrow \sim R\) & \(R\wedge \sim G\) \tabularnewline\hline
        &&&&&&&&\tabularnewline\hline
        &&&&&&&&\tabularnewline\hline
        &&&&&&&&\tabularnewline\hline
        &&&&&&&&\tabularnewline\hline        
    \end{tabular}
    \end{center}
\end{practice}

\clearpage

\begin{practice}{True or False: The Return}
    In practice problem \ref{GRimages} for which images are these true?
    \begin{enumerate}
    \setlength{\itemsep}{10pt}
        \item If ``there is a zombie,'' then ``the grass is wet.'' \hrulefill
        \item If ``there is a zombie,'' then ``it is raining.'' \hrulefill
        \item If ``it is raining,'' then ``the grass is green.'' \hrulefill
        \item If ``the grass is wet,'' then ``the grass is green.'' \hrulefill
    \end{enumerate}
\end{practice}



\begin{expository}{Conditionals and Their Relations}
    Starting with a conditional \(P \rightarrow Q\) we can write the following:
    \begin{multicols}{2}
        \begin{enumerate}
            \item \emph{Converse}\index{converse}: \(Q\rightarrow P\)
            \item \emph{Inverse}\index{inverse}: \(\sim P\rightarrow \sim Q\)
            \item \emph{Contrapositive}\index{contrapositive}: \(\sim Q\rightarrow \sim P\)
            \item \emph{Negation}\index{negation}: \(P\wedge \sim Q\)
        \end{enumerate}
    \end{multicols}
    Where \((P\rightarrow Q) \equiv (\sim Q \rightarrow \sim P)\), 
    \((Q\rightarrow P) \equiv (\sim P \rightarrow \sim Q)\), 
    and \(\sim (P\rightarrow Q) \equiv (P \wedge \sim Q)\).
    
    Further we can say \(P\rightarrow t\) and \(c\rightarrow Q\) are \emph{vacuously true}\index{vacuously true} because their negations are always false:
        \[\sim (P\rightarrow t) \equiv (P \wedge c) \equiv c\]
    and
        \[\sim (c\rightarrow Q) \equiv (c \wedge \sim Q) \equiv c.\]
\end{expository}

\setcounter{n}{1}
\begin{practice}{Making Arguments}
    If we assume it is true that ``If I carry an umbrella, then zombies dive bomb me,'' then which pictures should we look at? \answerbox
    \begin{center}
        \begin{tikzpicture}
            \foreach \p/\x in {PlaneU/0,U/4,Plane/8,None/12}{
                \pic[scale=0.5,transform shape] at (\x,0) {\p};
                \node at (\x-1,2.5) {(\then)};
                \stepcounter{n};
            }
        \end{tikzpicture}
    \end{center}
    What if we also assume that it is true that ``I am carrying an umbrella,''
    which picture(s) do we look at then \answerbox, are there zombies dive bombing me \answerbox?\vspace{\baselineskip}

    \noindent%
    What if we also assume that it is true that ``zombies dive bomb me,'' which picture(s) do we look at then \answerbox, am I carrying an umbrella \answerbox?
    
\end{practice}

\clearpage

\begin{expository}{More Rules and Arguments}
    Given statements \(P\) and \(Q\) the following are some common valid and invalid argument forms.
    \vspace{\baselineskip}

    \begin{tabular}{*{4}{p{0.22\textwidth}}}
    \emph{Modus Ponens}\index{modus ponens}:
        \[
            \begin{array}{ll}
                P\rightarrow Q  \\
                P \\\hline
                \therefore\ Q
            \end{array}
        \] &
    \emph{Modus Tollens}\index{modus tollens}:
        \[
            \begin{array}{ll}
                P\rightarrow Q  \\
                \sim Q \\\hline
                \therefore\ \sim P
            \end{array}
        \] &
    \emph{Converse Error}\index{converse error}:
        \[
            \begin{array}{ll}
                P\rightarrow Q  \\
                Q \\\hline
                \therefore\ P
            \end{array}
        \] &
    \emph{Inverse Error}\index{inverse error}:
        \[
            \begin{array}{ll}
                P\rightarrow Q  \\
                \sim P \\\hline
                \therefore\ \sim Q
            \end{array}
        \] \tabularnewline
    \emph{Disjunctive\newline Syllogism}\index{disjunctive syllogism}:
        \[
            \begin{array}{ll}
                P\vee Q  \\
                \sim P \\\hline
                \therefore\ Q
            \end{array}
        \] &
    \emph{Law of\newline Simplification}\index{law of simplification}:
        \[
            \begin{array}{ll}
                P\wedge Q  \\ \hline
                \therefore\ P\ \text{and}\ Q
            \end{array}
        \] &
    \emph{Transitivity}\index{transitivity}:
        \[
            \begin{array}{ll}
                P\rightarrow Q  \\
                Q\rightarrow R \\\hline
                \therefore\ P\rightarrow R
            \end{array}
        \] &
    \emph{Contradiction}\index{contradiction}:
        \[
            \begin{array}{ll}
                P \wedge \sim Q\rightarrow c  \\\hline
                \therefore\ P\rightarrow Q
            \end{array}
        \] 
    \end{tabular}
\end{expository}